\chapter{INTRODUCTION}

\section{Background}

Coastal and estuarine environments stand among the most dynamic, complex, and
hydrologically sensitive regions on Earth. Governed by the intricate interplay of
tidal forcing from the ocean, continuous riverine discharge from upstream catchments,
and the complex geometry of local bathymetry, free-surface water flows in these transitional zones dictate a broad spectrum of
critical environmental and societal outcomes. These include flood inundation dynamics, sediment
transport mechanics, salinity intrusion into freshwater aquifers, ecological habitat sustainability, and
the safety of maritime navigation \cite{Sottolichio2000, Jalón-Rojas2018}. The ability to
accurately predict water surface elevation and velocity fields across an
estuary is therefore of fundamental and growing importance to coastal engineers,
environmental managers, urban planners, and climate scientists alike.

The global threat of climate change has rendered accurate hydrodynamic
prediction more urgent than ever before. Projections from the
Intergovernmental Panel on Climate Change (IPCC) and numerous subsequent studies indicate that rising mean sea
levels, combined with more frequent and intense storm surges and shifting river
discharge regimes, will substantially increase the frequency, spatial extent, and severity
of coastal flooding worldwide \cite{IPCC2021, Vousdoukas2018, Alfieri2017}.
In macro-tidal estuaries---where natural tidal ranges may exceed several metres
and the non-linear coupling between oceanic tides and extreme riverine floods
generates highly complex and unpredictable inundation patterns---even marginal inaccuracies in
predicted water levels can have catastrophic consequences for human
populations, critical infrastructure, and fragile coastal ecosystems. Consequently, robust flood risk assessment and real-time forecasting systems are imperative.

Historically, physics-based numerical hydrodynamic models have served as the
foundational standard for simulating complex estuarine flows. Industry-standard solvers such as Deltares'
Delft3D Flexible Mesh (Delft3D FM) \cite{Delft3DFM2011}, TELEMAC-2D \cite{Hervouet2000},
and HEC-RAS 2D \cite{Quirogaetal2020} discretise the depth-averaged
Shallow Water Equations (SWE) on high-resolution structured or unstructured computational
meshes. These models simulate tidal propagation, wave dynamics, and current interactions with exceptionally high physical fidelity. When
carefully calibrated against empirical field observations, such as tide gauge records and Acoustic Doppler Current Profiler (ADCP) measurements, these models reliably
reproduce critical hydrodynamic phenomena, including tidal constituents, tidal asymmetry, and the non-linear
interactions between river and ocean forcing. However, their extraordinary computational cost presents a major operational limitation. This cost is primarily driven by the Courant--Friedrichs--Lewy (CFL) condition---a strict numerical stability constraint that restricts the simulation Courant number, forcing the solver to take extremely small time steps on high-resolution meshes. Consequently, running these simulations often requires hours or even days of
wall-clock time per tidal cycle on modern multi-core high-performance computing (HPC)
clusters. This computational burden constitutes a fundamental bottleneck for applications that
demand rapid iteration, such as large ensemble runs for probabilistic forecasting, real-time operational flood warning systems, comprehensive uncertainty
quantification, or extensive large-scale climate-impact scenario analysis
\cite{Bates2010, Brunton2020}.

In response to these computational challenges, the past decade has witnessed a rapid proliferation of machine learning (ML)
and deep learning (DL) methods across the earth and environmental sciences
\cite{LeCun2015, Reichstein2019, Shen2018}. Various architectures, including traditional Artificial Neural Networks (ANNs),
Long Short-Term Memory (LSTM) networks for sequential data, Convolutional Neural Networks (CNNs) for spatial grids, and,
most recently, Graph Neural Networks (GNNs) for unstructured spatial data, have been successfully applied to problems
ranging from rainfall-runoff modelling and urban flood mapping to global ocean state
forecasting \cite{Sit2020, Bentivoglio2023}. Once trained offline, these
data-driven surrogate models can generate predictions in milliseconds,
offering computational speed-up factors of several orders of magnitude relative to
classical physics-based solvers. Nevertheless, purely data-driven models suffer from a critical flaw: they are fundamentally agnostic to the governing physical laws of fluid mechanics. Consequently, they are prone to producing non-physical, highly erroneous predictions when deployed in out-of-distribution scenarios or extreme conditions not well-represented in their training data \cite{Karniadakis2021, Cuomo2022}.

A principled resolution to this dichotomy between computational speed and physical fidelity is provided by
the paradigm of \textbf{Physics-Informed Neural Networks (PINNs)}, introduced by Raissi,
Perdikaris, and Karniadakis \cite{Raissi2019}. In the PINN framework,
the residuals of the governing partial differential equations (PDEs)---in this case, the Shallow Water Equations---are
evaluated via automatic differentiation and embedded directly as penalty terms in
the neural network's training loss function. This innovative approach forces the network to honour both
the observed empirical data and the underlying physical conservation laws simultaneously. The result is a predictive model that yields solutions which are physically consistent, highly accurate, and robust even in data-sparse
regions or under extrapolated boundary conditions. Extensions of PINNs to solve the one- and two-dimensional Shallow
Water Equations have been successfully demonstrated in recent exploratory studies
\cite{Tian2025, Yin2024, Dazzi2024}, confirming their immense potential for
hydrodynamic surrogate modelling.

Despite the success of PINNs, a significant architectural challenge arises when applying them to the highly irregular geometry
and complex bathymetry of real-world estuaries. Standard PINN architectures, which typically operate on simple continuous coordinates or regular Cartesian grids, cannot
adequately represent the complex, unstructured finite-volume meshes commonly used by state-of-the-art
numerical solvers like Delft3D FM. To overcome this, \textbf{Graph Neural Networks (GNNs)}
\cite{Scarselli2009, Gilmer2017} are employed. GNNs address this limitation naturally by
treating the computational mesh as a graph structure---where mesh cells function as nodes and shared faces act as edges. This allows the network to perform iterative message passing, exchanging spatial information locally in a
topology-aware manner that perfectly mirrors the numerical flux exchanges in finite-volume methods. Recent advances, such as the Data-Guided Finite Volume-Informed Neural Network (FVM-PINN) introduced by Liu \cite{Liu2026}, explicitly highlight that incorporating finite-volume physics and sparse data guidance is critical to preventing standard strong-form PINNs from failing on complex unstructured grids. Merging the physical rigour and PDE-awareness of PINNs with the
geometric flexibility and spatial scalability of GNNs yields the highly advanced \textbf{Physics-Informed Graph
Neural Network (PI-GNN)} framework \cite{Taghizadeh2025}. This synthesis forms the core
methodological foundation of this thesis, aiming to provide a fast, accurate, and physically consistent surrogate for estuarine hydrodynamics.

\section{Context}

The focal point of this study is the \textbf{Gironde estuary}, located in south-western
France. Formed by the confluence of the Garonne and Dordogne rivers near the Bec d'Ambès and flowing into the Atlantic Ocean at Pointe de Grave, the Gironde is the largest macro-tidal estuary in Western Europe, covering a surface area of approximately 635 square kilometres.
Stretching roughly 75\,km in length, the estuary is characterised by extremely dynamic hydrodynamics. It experiences mean spring tidal ranges exceeding 5\,m at its mouth, which amplify as they propagate upstream, leading to pronounced
non-linear tidal distortion and the generation of significant shallow-water tidal constituents (e.g., M4, M6). Additionally, the estuary receives a highly variable seasonal combined riverine discharge from the Garonne and Dordogne rivers, ranging from less than 200\,m$^3$/s during dry summer months to overwhelming peak flows exceeding 5\,000\,m$^3$/s during winter and spring flood events \cite{Sottolichio2000,
Jalón-Rojas2018}. 

The Gironde estuary is a region of immense economic and ecological importance. It supports major shipping and maritime traffic navigating to the inland
port of Bordeaux, sustains extensive viticultural and agricultural lowlands along its banks, and provides critical habitats for numerous marine and avian species. Furthermore, it borders the greater Bordeaux metropolitan area, home to more than 800\,000 inhabitants. The combination of dense population, vital infrastructure, and extreme tidal-fluvial interactions makes the region exceptionally vulnerable to complex flooding scenarios, necessitating highly accurate and responsive hydrodynamic monitoring and prediction systems.

To capture the dynamics of this complex system, a high-resolution \textbf{Delft3D Flexible Mesh (Delft3D FM)} numerical model of
the Gironde estuary was independently developed, calibrated, and validated specifically for this
research. This custom-built model encompasses a highly detailed domain of 36\,271 unstructured finite-volume cells
and 53\,224 internal faces, serving as the primary source of high-fidelity ground-truth data for this
thesis. The numerical model comprehensively resolves the full estuarine domain---extending from the Atlantic
open boundary into the upstream tidal limits of both the Garonne and
Dordogne rivers. It has been rigorously calibrated and validated against historical observed water surface elevation
and velocity data. The simulation output, precisely recorded at 30-minute intervals over
a continuous 265-hour tidal cycle, provides a comprehensive, multidimensional dataset of water surface
elevation ($\eta$), depth-averaged longitudinal velocity ($u$), and lateral velocity ($v$). This dataset forms the foundation against which the proposed PI-GNN surrogate model is trained, constrained, and ultimately tested.

The specific computational challenge that motivates this research is the critical need for
a robust \textbf{parametric surrogate model}. Unlike simple spatial interpolators, a parametric surrogate must accept not only the fixed
spatio-temporal coordinates ($x, y, t$) and static bathymetric depth ($z$) as inputs, but also dynamic, time-varying boundary conditions. Specifically, the model must be conditioned on the
tidal amplitude at the open ocean boundary ($H_\mathrm{ocean}$) and the time-varying
river discharges from the Garonne ($Q_\mathrm{Garonne}$) and Dordogne
($Q_\mathrm{Dordogne}$) rivers. A successfully trained parametric
surrogate can then instantly predict the full hydrodynamic state
$[\eta,\, u,\, v]$ at any spatial node within the estuary, at any given time, for \emph{any} arbitrary
combination of these boundary forcings. Crucially, it accomplishes this without requiring a costly re-simulation of the underlying physics-based numerical model. This transformative capability is
particularly valuable for enabling real-time operational flood forecasting, conducting ensemble-based
probabilistic risk assessments, and exploring diverse, long-term climate-change scenarios with unprecedented computational efficiency.

However, developing such a surrogate for tidal environments presents a profound technical challenge known as
\textbf{spectral bias}. In the context of deep learning, standard multilayer perceptron (MLP) networks
exhibit a strong propensity to learn low-frequency, slowly varying components of a target signal rapidly, while fundamentally
struggling to resolve high-frequency, oscillatory components \cite{Tancik2020}. Because estuarine hydrodynamics are fundamentally dominated by high-frequency tidal oscillations (such as the dominant semi-diurnal $\sim$12\,h and quarter-diurnal $\sim$6\,h periods), a standard neural network will fail to capture the complex tidal waves, resulting in overly smoothed, inaccurate predictions. This thesis directly addresses the critical issue of spectral bias
through the novel implementation of \textbf{Random Fourier Feature (RFF) embeddings}. By projecting the input coordinates into a higher-dimensional periodic space with
independently tuned temporal and spatial bandwidths, the RFF encoding effectively allows the PI-GNN to perceive and accurately reconstruct the complex, high-frequency tidal dynamics inherent to the estuarine system.

\section{Aims and Objectives}

The overarching \textbf{aim} of this Master's thesis is to conceptualise, develop, train, and
rigorously validate a state-of-the-art parametric Physics-Informed Graph Neural Network
(PI-GNN) surrogate model. This model is designed to accurately and instantaneously emulate the complex 2D tidal hydrodynamics of the
macro-tidal Gironde estuary. This surrogate learns directly from high-fidelity simulation data generated by the newly developed and calibrated Delft3D FM model, while being explicitly constrained by the governing physical laws of fluid motion. The resulting surrogate must map a comprehensive seven-dimensional input vector
$[\,t,\, x,\, y,\, z,\, H_\mathrm{ocean},\, Q_\mathrm{Garonne},\,
Q_\mathrm{Dordogne}\,]$ to predict the full hydrodynamic state variables
$[\,\eta,\, u,\, v\,]$, perfectly conditioned on dynamic boundary forcings, thereby bypassing the need for computationally prohibitive numerical simulations.

To achieve this overarching aim, the study pursues the following specific \textbf{objectives}:

\begin{enumerate}
    \item \textbf{Numerical Modelling and Data Curation:} To independently develop and calibrate a high-fidelity Delft3D FM numerical model of the Gironde
    estuary. Subsequently, to meticulously extract, process, and normalize a massive training
    dataset from its simulations. This includes processing the complex unstructured mesh geometry, extracting cell-centred
    bathymetric depths, and preparing the time-series data of water surface elevation and depth-averaged
    velocity fields across the entire 265-hour simulation period, ensuring optimal formatting for graph-based deep learning.

    \item \textbf{Architecture Design with Spectral Bias Mitigation:} To design and engineer a robust PI-GNN architecture that explicitly overcomes spectral bias. This involves integrating Random Fourier
    Feature (RFF) embeddings---utilising highly specific and decoupled temporal ($\sigma_t = 30$)
    and spatial ($\sigma_s = 1$) bandwidths---with a deep, six-layer
    message-passing Multi-Layer Perceptron (comprising 512 hidden units per layer with SiLU activations). This network will operate directly on the
    unstructured finite-volume cell graph derived from the Gironde computational mesh.

    \item \textbf{Physics-Informed Loss Formulation:} To mathematically formulate and programmatically implement a comprehensive, autograd-based physics-informed
    loss function. This function must explicitly encode the complete 2D Shallow Water Equations, encompassing mass
    conservation, x-momentum, and y-momentum. Crucially, it must incorporate complex bed-slope corrections based on the estuary's bathymetry. This physics loss will be optimized alongside traditional data fidelity (Mean Squared Error), boundary condition, and initial condition
    terms, managed via an adaptive, multi-stage
    curriculum training schedule to ensure stable convergence.

    \item \textbf{Model Validation and Computational Acceleration:} To comprehensively validate the predictive accuracy and physical consistency of the fully trained PI-GNN surrogate on a temporally
    held-out test dataset (representing the final 20\% of the simulation, or approximately 55
    hours of entirely unseen tidal forcing). Furthermore, to explicitly quantify the dramatic reduction in simulation inference time achieved by the surrogate relative to the traditional numerical solver. Performance will be quantified using standard hydrological metrics, including Root Mean Square Error (RMSE),
    Nash--Sutcliffe Efficiency (NSE), and the coefficient of
    determination ($R^2$), evaluated rigorously at eight designated observation stations and
    five strategically located interior nodes across the estuary.

    \item \textbf{Assessment of Parametric Generalisation:} To critically assess the parametric generalisation capabilities of the
    surrogate model by evaluating its predictive accuracy and stability when subjected to novel combinations of
    tidal amplitudes and riverine discharges that were not encountered during the model's training phase. This objective aims to definitively demonstrate the model's operational utility as a high-speed digital twin for practical scenario
    analysis and flood forecasting.
\end{enumerate}

\section{Significance, Scope and Definitions}

\subsection*{Significance}

This research constitutes a significant advancement at the intersection of
computational coastal hydrology and applied scientific machine learning. It makes several original and impactful contributions to the field:

\begin{itemize}
    \item \textbf{Large-Scale Parametric Emulation:} It presents one of the first successful implementations of a fully \textbf{parametric PI-GNN
    surrogate model} for modelling 2D tidal hydrodynamics on a real, large-scale,
    highly dynamic macro-tidal estuary. The successful handling of an unstructured finite-volume mesh comprising
    more than 36\,000 cells demonstrates the scalability of graph-based methods for real-world coastal engineering applications.

    \item \textbf{Advanced Physics Integration:} It demonstrates the successful practical implementation of highly complex,
    \textbf{autograd-based 2D SWE physics residuals} within a graph-structured deep learning framework. By successfully encoding mass
    conservation and momentum equations, including critical bathymetric bed-slope
    corrections computed via Green--Gauss gradient reconstruction, it pushes the boundaries of how much physical realism can be embedded in neural networks for irregular domains.

    \item \textbf{Overcoming Spectral Bias in Tidal Modelling:} It introduces and validates a novel \textbf{dual-bandwidth Random Fourier Feature
    encoding} scheme that is specifically tailored to the distinct temporal (rapid tidal) and
    spatial (large-scale estuarine) length scales of coastal problems. This directly and effectively addresses
    the spectral bias phenomenon that typically causes standard PINNs to fail when modelling highly oscillatory tidal signals.

    \item \textbf{Transformative Computational Acceleration:} It definitively proves that complex, physics-constrained tidal hydrodynamics can be simulated in fractions of a second during inference. By drastically reducing the computational time compared to traditional CPU-bound numerical solvers, it paves the way for real-time flood warning systems and large-scale ensemble forecasting.

    \item \textbf{Open-Source Methodology:} It provides a fully \textbf{reproducible and open-source training pipeline}, which has been thoroughly tested and optimized on GPU-accelerated computing platforms (such as Kaggle and local CUDA environments). This pipeline is highly adaptable and can be seamlessly applied to other complex estuarine or coastal basins worldwide, provided they have baseline calibrated numerical simulation datasets.
\end{itemize}

\subsection*{Scope}

To ensure focus and computational feasibility, the scope of this study is specifically limited to the following parameters:

\begin{itemize}
    \item \textbf{Physical Equations:} The modelling is restricted strictly to \textbf{depth-averaged, 2D free-surface hydrodynamics} as governed by the Shallow Water Equations. Complex three-dimensional flow structures, advanced turbulence closure models (e.g., $k-\epsilon$), and density stratification effects (baroclinic flows) are explicitly excluded from this analysis.

    \item \textbf{Geographical Domain:} The study is exclusively focused on the \textbf{Gironde estuary} in France. The simulations utilise data generated by the Delft3D FM numerical model developed as part of this thesis, applying a constant, uniform Manning's bed roughness coefficient of $n = 0.019$\,s\,m$^{-1/3}$ across the entire domain.

    \item \textbf{Target Variables:} The surrogate model is trained to predict \textbf{hydrodynamic state variables only}: specifically, water surface
    elevation ($\eta$), depth-averaged eastward velocity ($u$), and
    depth-averaged northward velocity ($v$). Phenomena such as sediment transport, morphological evolution, wave action, and salinity or pollutant intrusion are beyond the current scope.

    \item \textbf{Temporal Domain:} The analysis is based on a continuous \textbf{265-hour simulation period}, which captures multiple distinct tidal cycles. This dataset is rigorously split using an 80/20 chronological ratio, resulting in a training set of approximately 210 hours and a temporally distinct, held-out test set of approximately 55 hours for uncompromised validation.
\end{itemize}

\subsection*{Definitions}

To ensure clarity and precision, key technical terms utilized extensively throughout this thesis are defined below:

\begin{description}
    \item[Shallow Water Equations (SWE):] The depth-averaged, vertically
    integrated form of the Navier--Stokes equations. They govern free-surface
    fluid flows in scenarios where the horizontal length scale of the fluid motion greatly exceeds the vertical water
    depth, describing the fundamental conservation of mass and depth-averaged
    momentum in two horizontal spatial dimensions.

    \item[Physics-Informed Neural Network (PINN):] An advanced class of neural networks trained
    by minimizing a composite loss function. This function uniquely incorporates traditional data
    fidelity terms (to match observations) alongside PDE
    residual terms. These residual terms act as soft penalties for any violations of the underlying governing physical equations,
    computed efficiently using automatic differentiation within the neural network framework.

    \item[Graph Neural Network (GNN):] A highly specialized category of deep learning architectures designed to
    operate directly on arbitrarily structured graph data. GNNs function through iterative
    message passing and aggregation mechanisms between connected nodes. This allows the network to
    learn complex representations that explicitly respect the geometric topology and connectivity of the underlying
    graph---in this context, the unstructured computational mesh of the estuary.

    \item[PI-GNN (Physics-Informed Graph Neural Network):] A novel hybrid deep learning
    architecture that synergistically combines the rigorous physical constraint enforcement characteristic of
    PINNs with the unparalleled geometric flexibility and spatial awareness of GNNs. It is specifically designed for highly accurate,
    physics-consistent surrogate modelling on complex, irregular computational
    meshes.

    \item[Surrogate Model:] A highly optimized, computationally inexpensive mathematical or algorithmic approximation
    of a complex, high-fidelity numerical simulator. By learning complex input-output mappings from pre-computed simulation data, a surrogate is
    capable of producing highly accurate predictions in a tiny fraction of the time required
    by the original, computationally heavy physics-based solver.

    \item[Random Fourier Feature (RFF) Embedding:] A sophisticated mathematical technique used to map low-dimensional input coordinates (such as time $t$ or spatial coordinates $x, y$)
    into a much higher-dimensional feature space using specifically tuned, random sinusoidal
    projections. This embedding is crucial for enabling deep neural networks to accurately learn high-frequency
    functions, effectively overcoming the inherent limitation known as spectral bias \cite{Tancik2020}.

    \item[Parametric Surrogate:] An advanced type of surrogate model that is designed to condition its
    predictions dynamically based on explicit, variable boundary conditions or forcing parameters (in this study, ocean
    tidal amplitude and various upstream riverine discharges). This parametric capability enables rapid, on-the-fly prediction
    across a virtually infinite family of different environmental scenarios without the need for model retraining.

    \item[Curriculum Learning:] A strategic, phased machine learning training paradigm. In curriculum learning, the
    difficulty, complexity, or composition of the learning task is deliberately and progressively increased
    during the training process. In the context of PINNs, this often involves initially training the network to fit only the empirical data without physics constraints, and then gradually and systematically introducing the complex PDE penalties as the model's accuracy improves \cite{Bengio2009}.

    \item[Manning's Roughness Coefficient ($n$):] An essential empirical parameter utilised
    in the SWE friction term (specifically, the Manning--Strickler formulation). It represents
    the physical resistance to fluid flow offered by the roughness of the channel bed. For this study, it is defined as
    a uniform constant of $n = 0.019$\,s\,m$^{-1/3}$ across the entire estuarine domain.

    \item[Nash--Sutcliffe Efficiency (NSE):] A robust, dimensionless statistical metric
    (ranging from $-\infty$ to 1) that is universally employed in hydrology to rigorously quantify the predictive skill of a model. An NSE value of $1$ denotes a perfect match between model predictions and observed data, an NSE of $0$ indicates that the model predictions are only as accurate as the mean of the observed data, and an
    NSE $< 0$ indicates that the model is fundamentally flawed and performs worse than simply using the observed mean.
\end{description}

\section{Thesis Outline}

The remainder of this thesis is structured to provide a comprehensive and logical progression of the research, organized as follows:

\textbf{Chapter~2 (Literature Review)} provides an exhaustive and critical review of the existing body of
knowledge. It systematically explores: (i)~the foundational principles of traditional numerical hydrodynamic modelling and its historical
application to complex estuarine systems; (ii)~the rapid theoretical development and practical application
of Physics-Informed Neural Networks (PINNs) for solving complex PDEs in fluid mechanics; (iii)~the emergence of Graph Neural Networks (GNNs) for analyzing unstructured-mesh hydrodynamics and
their critical role in accelerating traditional numerical solvers; and (iv)~state-of-the-art surrogate
modelling strategies for real-time coastal flood mapping and tidal forecasting.
The chapter concludes by synthesizing a clear summary of the identified critical research gaps that directly
motivate the objectives of the present study.

\textbf{Chapter~3 (Materials and Methodology)} provides a deeply detailed exposition of the research materials and the proposed methodologies. It begins by describing the complex geophysical and hydrodynamic characteristics of the study area
(the Gironde estuary). It then details the independent development and configuration of the Delft3D FM numerical model, the rigorous
calibration procedure applied against field observations, and the precise nature of the simulation dataset generated for neural network training and
validation. Following this, the chapter meticulously details the architecture of the proposed PI-GNN surrogate. This includes a deep dive into
the Random Fourier Feature encoding strategy, the design of the message-passing graph neural network structure, the complex mathematical formulation of the
autograd-based physics-informed loss function, the rationale behind the multi-stage curriculum training strategy,
and the specifics of the GPU-accelerated computational implementation.

\textbf{Chapter~4 (Results)} presents the comprehensive findings of the model training and validation phases. It includes detailed analyses of the training convergence curves and explicitly reports
quantitative performance metrics (including RMSE, NSE, and $R^2$) evaluated at specific observation
stations and across interior nodes for both the training phase and the critical held-out test
period. Furthermore, it presents vivid spatial comparisons of water-level and velocity fields, alongside descriptions of the animated
tidal simulation outputs generated by the fully trained PI-GNN surrogate.

\textbf{Chapter~5 (Analysis and Discussion)} provides a critical and in-depth analysis of the results. It interprets the
PI-GNN surrogate's performance explicitly in relation to the research gaps identified in Chapter 2. It rigorously
examines the individual contributions and impacts of the various components of the physics loss terms, compares the advanced PI-GNN model
against baseline purely data-driven models to highlight the value of physics integration, and extensively discusses the broader real-world implications of the research for operational
tidal forecasting systems and long-term climate-change impact assessments.

\textbf{Chapter~6 (Conclusions)} serves as the final synthesis of the thesis. It concisely summarises the principal findings of
the entire study, clearly restates the original contributions made to the scientific community, candidly acknowledges the
inherent limitations and boundaries of the present work, and provides thoughtful recommendations and clear directions for future research endeavours. These recommendations include exploring multi-fidelity training paradigms, investigating domain adaptation techniques to apply the model to other estuaries globally,
and extending the current architecture to encompass coupled hydrodynamic and sediment transport modelling.
